\begin{figure}[!htbp]
		\centering
		\begin{tikzpicture}[
			trave/.style={line width=2.0pt, line cap=round,
				line join=round, black},
			terra/.style={line width=0.8pt, black},
			nodoV/.style={circle, fill=white, draw=black,
				inner sep=0pt, minimum size=6pt, line width=0.9pt},
			]
			
			\def\Lp{2.4}   % unità di lunghezza (= ell): altezza = luce/2
			\def\dx{6.5}   % offset orizzontale tra i due pannelli
			
			%%═══════════════════════════════════════════════════════════════
			%% PANNELLO (a) — PORTALE A TRE CERNIERE
			%% A=(0,0), B=(0,L), C=(L,L), D=(2L,L), E=(2L,0)
			%%═══════════════════════════════════════════════════════════════
			\begin{scope}[xshift=0cm]
				
				%% ── Coordinate ─────────────────────────────────────────
				\coordinate (A) at (0,       0   );
				\coordinate (B) at (0,       \Lp );
				\coordinate (C) at (\Lp,     \Lp );
				\coordinate (D) at (2*\Lp,   \Lp );
				\coordinate (E) at (2*\Lp,   0   );
				
				%% ── Corpi: A-B-C e C-D-E ──────────────────────────────
				\draw[trave] (A) -- (B) -- (C);
				\draw[trave] (C) -- (D) -- (E);
				
				%% ── Cerniera esterna in A ──────────────────────────────
				\draw[terra] (A) -- ++(-0.25,-0.42) -- ++(0.50,0) -- cycle;
				\draw[terra] ($(A)+(-0.34,-0.42)$) -- ($(A)+(0.34,-0.42)$);	\foreach \x in {-0.28,-0.14,0,0.14,0.28}
				\draw[very thin,black] ($(A)+(\x,-0.42)$) -- ++(-0.10,-0.13);
				\node[nodoV] at (A) {};
				
				%% ── Cerniera esterna in E ──────────────────────────────
				\draw[terra] (E) -- ++(-0.25,-0.42) -- ++(0.50,0) -- cycle;
				\draw[terra] ($(E)+(-0.34,-0.42)$) -- ($(E)+(0.34,-0.42)$);	\foreach \x in {-0.28,-0.14,0,0.14,0.28}
				\draw[very thin,black] ($(E)+(\x,-0.42)$) -- ++(-0.10,-0.13);
				\node[nodoV] at (E) {};
				
				%% ── Cerniera interna in C (chiave) ─────────────────────
				\node[nodoV] at (C) {};
				
				%% ── Etichette dei nodi ─────────────────────────────────
				\node[left,  font=\small] at (A) {$A$};
				\node[above left,  font=\small] at (B) {$B$};
				\node[above,       font=\small] at ($(C)+(0,0.14)$) {$C$};
				\node[above right, font=\small] at (D) {$D$};
				\node[right, font=\small] at (E) {$E$};
				
				%% ── Etichette dei corpi ────────────────────────────────
				\node[font=\small] at ({0.60},{0.55*\Lp}) {$\mathcal{C}_1$};
				\node[font=\small] at ({2.0*\Lp-0.60},{0.55*\Lp}) {$\mathcal{C}_2$};
				
				%% ── Didascalia del pannello ────────────────────────────
				\node[font=\small] at (\Lp,0)
				{(a)};
				
			\end{scope}
			
			%%═══════════════════════════════════════════════════════════════
			%% PANNELLO (b) — PORTALE DOPPIAMENTE INCASTRATO
			%%═══════════════════════════════════════════════════════════════
			\begin{scope}[xshift=\dx cm]
				
				%% ── Coordinate ─────────────────────────────────────────
				\coordinate (A2) at (0,       0   );
				\coordinate (B2) at (0,       \Lp );
				\coordinate (C2) at (\Lp,     \Lp );
				\coordinate (D2) at (2*\Lp,   \Lp );
				\coordinate (E2) at (2*\Lp,   0   );
				
				%% ── Corpo unico: A-B-C-D-E (traversa continua) ────────
				\draw[trave] (A2) -- (B2) -- (D2) -- (E2);
				
				%% ── Incastro in A: parete orizzontale a y=0 sotto, hatching giù
				\draw[terra] ($(A2)+(-0.40, 0)$) -- ($(A2)+(0.40, 0)$);
				\foreach \x in {-0.34,-0.20,-0.06,0.08,0.22,0.36}
				\draw[very thin,black] ($(A2)+(\x,0)$) -- ++(-0.10,-0.13);
				
				%% ── Incastro in E: parete orizzontale a y=0 sotto, hatching giù
				\draw[terra] ($(E2)+(-0.40, 0)$) -- ($(E2)+(0.40, 0)$);
				\foreach \x in {-0.34,-0.20,-0.06,0.08,0.22,0.36}
				\draw[very thin,black] ($(E2)+(\x,0)$) -- ++(-0.10,-0.13);
				
				%% ── Etichette dei nodi ─────────────────────────────────
				\node[above left,  font=\small] at (A2) {$A$};
				\node[above left,  font=\small] at (B2) {$B$};
				%\node[above,       font=\small] at ($(C2)+(0,0.14)$) {$C$};
				\node[above right, font=\small] at (D2) {$D$};
				\node[above right, font=\small] at (E2) {$E$};
				
				%% ── Etichetta del corpo unico ──────────────────────────
				\node[font=\small] at (\Lp,{\Lp+0.3}) {$\mathcal{C}$};
				
				%% ── Didascalia del pannello ────────────────────────────
				\node[font=\small] at (\Lp,0)
				{(b)};
				
			\end{scope}
			
		\end{tikzpicture}
		\caption{(a) Portale a tre cerniere e (b) e portale	incastrato.}
		\label{fig:portali_confronto}
	\end{figure}
